% !TEX program = xelatex
% Самодостаточный конспект. Компилировать XeLaTeX/LuaLaTeX (tectonic 04-soprjazhennye-operatory.tex).
\documentclass[11pt,a4paper]{article}
\newcommand{\coursename}{Специальные разделы высшей математики}
\newcommand{\rightheader}{§4. Сопряжённые и самосопряжённые операторы}
\newcommand{\doctitle}{Сопряжённые и самосопряжённые операторы}
% ==== Общая преамбула конспектов (собирается в каждый .tex как самодостаточная) ====
% Перед этим файлом должны быть определены: \documentclass и макросы
%   \coursename  — название курса (левый колонтитул, подзаголовок),
%   \rightheader — правый колонтитул (напр. «§2. Рациональные дроби»),
%   \doctitle    — заголовок раздела на титуле.
\usepackage{fontspec}
\setmainfont{cmunrm.otf}[
  BoldFont       = cmunbx.otf,
  ItalicFont     = cmunti.otf,
  BoldItalicFont = cmunbi.otf]
\usepackage{polyglossia}
\setdefaultlanguage{russian}
\setotherlanguage{english}

\usepackage{amsmath,amssymb,amsthm,mathtools}
\usepackage[a4paper,margin=2.2cm]{geometry}
\usepackage{enumitem}
\usepackage{graphicx}
\usepackage{array}
\usepackage{xcolor}
\definecolor{accent}{HTML}{1F6FEB}
\definecolor{rulecol}{HTML}{D0D7DE}
\usepackage[colorlinks=true,linkcolor=accent,urlcolor=accent,citecolor=accent]{hyperref}
\usepackage{titlesec}
\usepackage{fancyhdr}

% ----- Колонтитулы -----
\pagestyle{fancy}
\fancyhf{}
\fancyhead[L]{\small\itshape \coursename}
\fancyhead[R]{\small\itshape \rightheader}
\fancyfoot[C]{\thepage}
\renewcommand{\headrulewidth}{0.4pt}
\renewcommand{\headrule}{\hbox to\headwidth{\color{rulecol}\leaders\hrule height \headrulewidth\hfill}}

% ----- Заголовки -----
\titleformat{\section}{\large\bfseries\color{accent}}{\thesection.}{0.6em}{}
\titleformat{\subsection}{\bfseries}{\thesubsection.}{0.5em}{}

% ----- Теоремные окружения (стиль конспекта) -----
\theoremstyle{definition}
\newtheorem{definition}{Определение}[section]
\newtheorem{example}[definition]{Пример}
\newtheorem{remark}[definition]{Замечание}
\theoremstyle{plain}
\newtheorem{theorem}[definition]{Теорема}
\newtheorem{lemma}[definition]{Лемма}
\newtheorem{corollary}[definition]{Следствие}
\newtheorem{property}[definition]{Свойство}
\newtheorem{criterion}[definition]{Критерий}
\renewcommand{\proofname}{Доказательство}

% ----- Операторы и сокращения -----
\DeclareMathOperator{\tg}{tg}
\DeclareMathOperator{\ctg}{ctg}
\DeclareMathOperator{\arctg}{arctg}
\DeclareMathOperator{\arcctg}{arcctg}
\DeclareMathOperator{\sh}{sh}
\DeclareMathOperator{\ch}{ch}
\DeclareMathOperator{\Th}{th}
\DeclareMathOperator{\cth}{cth}
\DeclareMathOperator{\Const}{const}
\DeclareMathOperator{\sign}{sign}
\DeclareMathOperator{\rang}{rang}
\DeclareMathOperator{\diam}{diam}
\DeclareMathOperator{\Span}{span}
\DeclareMathOperator{\Ker}{Ker}
\DeclareMathOperator{\Img}{Im}
\DeclareMathOperator{\tr}{tr}
\newcommand{\R}{\mathbb{R}}
\newcommand{\C}{\mathbb{C}}
\newcommand{\N}{\mathbb{N}}
\newcommand{\Z}{\mathbb{Z}}
\newcommand{\Q}{\mathbb{Q}}
\newcommand{\dx}{\,dx}
\newcommand{\dt}{\,dt}
\newcommand{\du}{\,du}
\newcommand{\eps}{\varepsilon}
\renewcommand{\Re}{\operatorname{Re}}
\renewcommand{\Im}{\operatorname{Im}}
\renewcommand{\le}{\leqslant}
\renewcommand{\ge}{\geqslant}
\newcommand{\scal}[2]{\left( #1,\,#2\right)}
\DeclarePairedDelimiter{\abs}{\lvert}{\rvert}
\DeclarePairedDelimiter{\norm}{\lVert}{\rVert}

\begin{document}
\begin{center}
  {\LARGE\bfseries \doctitle}\\[4pt]
  {\large\itshape \coursename}
\end{center}
\vspace{6pt}

\section{Сопряжённый оператор}

\begin{definition}
\emph{Сопряжённым} к оператору $\mathcal A$ в евклидовом пространстве называют оператор
$\mathcal A^*$, для которого
\[
  \scal{\mathcal Ax}{y}=\scal{x}{\mathcal A^*y}\qquad\forall x,y .
\]
\end{definition}

\begin{theorem}
$\mathcal A^*$ существует, единствен и линеен. В ортонормированном базисе его матрица —
транспонированная: $A^*=A^{T}$ (в унитарном случае $A^*=\overline{A}^{\,T}$).
\end{theorem}

\begin{proof}
В ОНБ $\scal{\mathcal Ax}{y}=(Ax)^{T}y=x^{T}A^{T}y=\scal{x}{A^{T}y}$, поэтому оператор с матрицей
$A^{T}$ удовлетворяет определению — существование. Единственность: если
$\scal x{\mathcal By}=\scal x{\mathcal Cy}$ для всех $x$, то $\scal x{(\mathcal B-\mathcal C)y}=0$
при всех $x$, в частности при $x=(\mathcal B-\mathcal C)y$, откуда $(\mathcal B-\mathcal C)y=0$ для
всех $y$. Линейность $\mathcal A^*$ следует из линейности по $y$ обеих частей определения. (Чисто
геометрически существование обеспечивает теорема Рисса о представлении линейного функционала
$x\mapsto\scal{\mathcal Ax}{y}$ вектором $\mathcal A^*y$.)
\end{proof}

\section{Алгебраические свойства сопряжения}

\begin{property}\leavevmode
\begin{enumerate}[label=\arabic*),leftmargin=*,nosep]
  \item $(\mathcal A+\mathcal B)^*=\mathcal A^*+\mathcal B^*$, $(\alpha\mathcal A)^*=\alpha\mathcal A^*$;
  \item $(\mathcal A\mathcal B)^*=\mathcal B^*\mathcal A^*$;
  \item инволютивность $(\mathcal A^*)^*=\mathcal A$;
  \item $\Img\mathcal A^*=(\Ker\mathcal A)^\perp$.
\end{enumerate}
\end{property}

\begin{proof}
2) $\scal{\mathcal A\mathcal Bx}{y}=\scal{\mathcal Bx}{\mathcal A^*y}=\scal{x}{\mathcal B^*\mathcal A^*y}$,
а по определению это $\scal x{(\mathcal A\mathcal B)^*y}$; в силу единственности
$(\mathcal A\mathcal B)^*=\mathcal B^*\mathcal A^*$.\;
4) $x\in\Ker\mathcal A\iff\mathcal Ax=0\iff\scal{\mathcal Ax}{y}=0\ \forall y
\iff\scal{x}{\mathcal A^*y}=0\ \forall y\iff x\perp\Img\mathcal A^*$. Значит
$\Ker\mathcal A=(\Img\mathcal A^*)^\perp$, откуда $\Img\mathcal A^*=(\Ker\mathcal A)^\perp$.
\end{proof}

\section{Самосопряжённые операторы}

\begin{definition}
Оператор \emph{самосопряжён} (симметрический), если $\mathcal A^*=\mathcal A$, т.е.\
$\scal{\mathcal Ax}{y}=\scal{x}{\mathcal Ay}$. В ОНБ это равносильно $A=A^{T}$ (симметрическая
матрица).
\end{definition}

\begin{remark}
Самосопряжённые операторы образуют подпространство в $\mathcal L(V,V)$: сумма и скалярное кратное
симметрических операторов симметричны.
\end{remark}

\section{Спектральные свойства самосопряжённых операторов}

\begin{theorem}
Все собственные значения самосопряжённого оператора вещественны.
\end{theorem}

\begin{proof}
Рассмотрим (при необходимости комплексифицировав) собственный вектор $x\ne0$, $\mathcal Ax=\lambda x$.
Тогда, используя эрмитово произведение и самосопряжённость,
\[
  \lambda\scal xx=\scal{\mathcal Ax}{x}=\scal{x}{\mathcal Ax}=\overline\lambda\,\scal xx .
\]
Так как $\scal xx>0$, получаем $\lambda=\overline\lambda$, т.е.\ $\lambda\in\R$.
\end{proof}

\begin{theorem}
Собственные векторы, отвечающие различным собственным значениям самосопряжённого оператора,
ортогональны.
\end{theorem}

\begin{proof}
Пусть $\mathcal Ax=\lambda x$, $\mathcal Ay=\mu y$, $\lambda\ne\mu$. Тогда
$\lambda\scal xy=\scal{\mathcal Ax}{y}=\scal{x}{\mathcal Ay}=\mu\scal xy$, откуда
$(\lambda-\mu)\scal xy=0$ и $\scal xy=0$.
\end{proof}

\section{Спектральная теорема}

\begin{lemma}
Если подпространство $L$ инвариантно относительно самосопряжённого $\mathcal A$
($\mathcal A L\subseteq L$), то и $L^\perp$ инвариантно.
\end{lemma}

\begin{proof}
Пусть $y\in L^\perp$. Для любого $x\in L$ имеем $\mathcal Ax\in L$, поэтому
$\scal{\mathcal Ay}{x}=\scal{y}{\mathcal Ax}=0$. Значит $\mathcal Ay\perp L$, т.е.\
$\mathcal Ay\in L^\perp$.
\end{proof}

\begin{theorem}[спектральная теорема]
Для самосопряжённого оператора в конечномерном евклидовом пространстве существует ортонормированный
базис из собственных векторов; в нём матрица диагональна.
\end{theorem}

\begin{proof}
Индукция по $\dim V$. При $\dim V=1$ утверждение очевидно. Шаг: характеристический многочлен имеет
вещественный корень $\lambda_1$ (собственные значения вещественны), ему отвечает единичный
собственный вектор $e_1$. Прямая $L=\Span(e_1)$ инвариантна, значит по лемме инвариантно
$L^\perp$ ($\dim L^\perp=\dim V-1$). Сужение $\mathcal A|_{L^\perp}$ самосопряжённо, по
предположению индукции в $L^\perp$ есть ОНБ из собственных векторов $e_2,\dots,e_n$. Вместе с $e_1$
они дают искомый ОНБ.
\end{proof}

\begin{remark}[геометрический смысл]
В этом базисе $\mathcal A=\sum_\lambda\lambda\,P_\lambda$, где $P_\lambda$ — ортогональный проектор
на $V_\lambda$: оператор «растягивает» взаимно ортогональные собственные направления с
коэффициентами $\lambda$.
\end{remark}

\end{document}
